\documentclass[11pt,a4paper]{article}
\usepackage[T1]{fontenc}
\usepackage[utf8]{inputenc}
\usepackage{lmodern}
\usepackage[margin=28mm]{geometry}
\usepackage{amsmath,amssymb,amsthm,mathtools}
\usepackage{microtype}
\usepackage[hidelinks]{hyperref}
\hypersetup{pdftitle={Small-charge rigidity for the quartic Hartree shape energy},pdfauthor={Anonymous},pdfsubject={Unrefereed mathematical research candidate},pdfkeywords={Hartree energy, shape optimisation, Coulomb interaction, small charge, ball rigidity, free boundary}}
\usepackage{enumitem}
\setlist{itemsep=3pt,topsep=5pt}
\newtheorem{theorem}{Theorem}[section]
\newtheorem{lemma}[theorem]{Lemma}
\newtheorem{proposition}[theorem]{Proposition}
\newtheorem{corollary}[theorem]{Corollary}
\theoremstyle{definition}
\newtheorem{definition}[theorem]{Definition}
\theoremstyle{remark}
\newtheorem{remark}[theorem]{Remark}
\DeclareMathOperator{\capop}{cap}
\DeclareMathOperator{\dist}{dist}
\DeclareMathOperator{\diver}{div}
\DeclareMathOperator{\tr}{tr}
\newcommand{\R}{\mathbb R}
\newcommand{\Sph}{\mathbb S^2}
\newcommand{\aesp}{\mathrm{ae}}
\newcommand{\qe}{\mathrm{qe}}
\newcommand{\dd}{\,\mathrm d}
\title{Small-charge rigidity for the quartic Hartree shape energy}
\author{Anonymous}
\date{9 October 2026}
\begin{document}
\maketitle
\begin{abstract}
We study the volume-constrained shape optimisation problem for a normalised Dirichlet state with repulsive quartic Coulomb interaction in three dimensions. We prove that, for every sufficiently small positive charge, the ball is the unique global minimiser up to translation and the natural Sobolev representative convention. The result uses the existence and all-minimiser near-sphericity theorem of Mazzoleni, Muratov and Ruffini. The remaining rigidity follows from an implicit-function argument for the physical boundary flux. We establish the material differentiability of the self-consistent state, identify its positive local branch with the actual ground state by an exact energy identity, and verify invertibility after removing volume and translation directions. A fixed-H\"older interpolation argument makes the local conclusion applicable to every global minimiser. We treat arbitrary finite-volume measurable supports as well as quasi-open domains and distinguish their equality conventions.
\end{abstract}

\section{Statement and relation to earlier work}\label{sec:intro}
Let $B=B_1(0)\subset\R^3$ and $V=|B|=4\pi/3$. For real functions define
\[
 D(f,g)=\iint_{\R^3\times\R^3}\frac{f(x)g(y)}{|x-y|}\dd x\dd y,
 \qquad
 \mathcal E_q(u)=\int_{\R^3}|\nabla u|^2\dd x+\frac q2D(u^2,u^2).
\]
Every state is normalised by $\|u\|_{L^2(\R^3)}=1$. For $u\in H^1(\R^3)$ its square belongs to $L^{6/5}$ by Sobolev interpolation, so the Coulomb term is finite by the Hardy--Littlewood--Sobolev inequality. All domains and states below are real and all charges in the optimisation assertion are positive.

We distinguish two standard domain conventions. If $\Omega$ is quasi-open, set
\[
 H^1_0(\Omega)=\{u\in H^1(\R^3):\widetilde u=0\text{ quasi-everywhere on }\R^3\setminus\Omega\},
 \qquad
 E_q(\Omega)=\inf_{\substack{u\in H^1_0(\Omega)\\\|u\|_2=1}}\mathcal E_q(u).
\]
Here $\widetilde u$ is a quasi-continuous representative. Capacity means Sobolev $H^1$ capacity; ``quasi-everywhere'' means outside a set of capacity zero. A set is quasi-open if it is open after adjoining sets of arbitrarily small capacity; equivalently it is a Sobolev positivity set up to capacity zero. The Newtonian capacity used in \cite[Section 2.3]{MMR2025} has the same null sets in three dimensions. On a fixed bounded set, multiplying a Newtonian admissible function by a fixed cutoff and using the Sobolev inequality controls its $H^1$ norm by its homogeneous gradient norm. The reverse comparison follows by dropping the $L^2$ term. Exhaustion and countable subadditivity extend the equivalence of null sets to arbitrary sets. We use only this equivalence, not numerical equality of the two capacities.

For an arbitrary Lebesgue measurable $\Omega$, the measurable-support relaxation is
\[
 X(\Omega)=\{u\in H^1(\R^3):u=0\text{ a.e. on }\R^3\setminus\Omega\},
 \qquad
 E_q^{\aesp}(\Omega)=\inf_{\substack{u\in X(\Omega)\\\|u\|_2=1}}\mathcal E_q(u).
\]
An infimum over an empty normalised class is $+\infty$. On a Lipschitz domain these definitions agree, by the zero-extension characterisation of $H^1_0$. On general quasi-open sets they need not agree: removing a set of positive capacity but zero volume can change $H^1_0$. The distinction is relevant to equality and is retained throughout the proof.

\begin{theorem}[Global small-charge rigidity]\label{thm:main}
There exists $q_0>0$ such that, for every $0<q<q_0$, the following hold.
\begin{enumerate}[label=(\roman*)]
\item For every quasi-open $\Omega\subset\R^3$ with $|\Omega|=V$,
\[
 E_q(\Omega)\geq E_q(B).
\]
Equality holds if and only if there is $z\in\R^3$ such that
$\capop\bigl(\Omega\mathbin\triangle(z+B)\bigr)=0$.
\item For every Lebesgue measurable $\Omega\subset\R^3$ with $|\Omega|=V$,
\[
 E_q^{\aesp}(\Omega)\geq E_q(B).
\]
Equality holds if and only if $\Omega=z+B$ almost everywhere for some $z\in\R^3$.
\end{enumerate}
No boundedness, connectedness, regularity or proximity assumption is imposed on the competitors.
\end{theorem}

Mazzoleni, Muratov and Ruffini introduced this shape problem in \cite{MMR2025}. Their Theorem 1.2 proves existence for small charge and shows that every minimiser, in its regular representative and modulo translation, has an arbitrarily small $C^{2,\alpha}$ radial graph. Their Remark 1.3 asks whether the ball is the unique minimiser. Theorem~\ref{thm:main} supplies that rigidity conclusion. We use their global compactness and free-boundary regularity theorem as an input; we do not reprove it.

The coupled interaction $D(u^2,u^2)$ differs from the domain-indicator Riesz energy in the earlier functional $\lambda_1(\Omega)+\varepsilon V_\alpha(\Omega)$. Small-repulsion ball optimality for the latter was established by Mazzoleni and Ruffini \cite{MazzoleniRuffini2021}. Here the repulsive density is selected jointly with the state, rather than fixed by the set. The quadratic interaction $D(u,u)$ studied in \cite{MMR2026quadratic} is also a different problem; none of its rigidity conclusions is assumed here.

Our local argument uses established implicit-function machinery for extremal domains. Pacard and Sicbaldi \cite{PacardSicbaldi2009} analyse the normalised eigenpair, the boundary operator and its translation kernel. Cavallina \cite{Cavallina2024} gives general symmetry and local uniqueness principles for nondegenerate critical shapes. In a different whole-space Hartree problem, Georgiev, Prinari and Visciglia \cite{GPV2012} already combine localisation of minimisers, removal of translations and implicit-function uniqueness to obtain radiality. Their energy includes a local attractive term and has no support-volume constraint or moving boundary. We do not claim these principles, or the spherical linear spectrum, as new. We verify the required differentiability and ground-state selection for the quartic Hartree shape functional and give the explicit specialisation and global transfer. This verification uses the first variation and the derivative of its physical boundary flux, and requires no quantified second-order expansion of the energy.

The structure is as follows. Section~\ref{sec:state} establishes the differentiable material state map and the boundary condition. Section~\ref{sec:local} proves local rigidity in one fixed H\"older topology. Section~\ref{sec:global} obtains uniform entrance of all global minimisers into that topology and proves Theorem~\ref{thm:main}, including the representative-sensitive equality assertions.


The parameterised branch theorem of Cavallina \cite[Theorem~I]{Cavallina2024families} is a further general precedent: its differentiable boundary map and bijective linearisation hypotheses yield local uniqueness. Those hypotheses must still be verified for the present coupled problem. For the same quartic functional, \cite[Theorem~1.1]{MMR2026screened} concerns nonexistence at large charge, not the small-charge rigidity proved here.

\section{The material ground state and its boundary flux}
\label{sec:state}

Two interfaces deserve particular care. Lemma~\ref{lem:selection} identifies the constructed positive branch with the full state infimum, so its flux describes the optimised energy. Later, Lemma~\ref{lem:entrance} places every global minimiser in one fixed H\"older neighbourhood before local uniqueness is applied.

Throughout this section, $B=B_1(0)\subset\mathbb R^3$ and
$0<\beta<1$ is fixed.  Functions in $H^1_0(\Omega)$ are extended by zero
when they occur in an integral over $\mathbb R^3$.  We use the notation
\[
 D(f,g)=\iint_{\mathbb R^3\times\mathbb R^3}
       \frac{f(x)g(y)}{|x-y|}\,dx\,dy,
 \qquad
 \mathcal E_q(u)=\int_{\mathbb R^3}|\nabla u|^2
                  +\frac q2D(u^2,u^2),
\]
and
\[
 E_q(\Omega)=\inf\{\mathcal E_q(u):u\in H^1_0(\Omega),
                              \|u\|_{L^2(\Omega)}=1\}.
\]
All states in this section are real-valued.

Choose $\chi\in C_c^\infty((1/2,3/2))$ that equals $1$ near $1$.
The bounded linear extension
\[
 (E\rho)(r\omega)=\chi(r)\rho(\omega)\omega,
 \qquad \rho\in C^{2,\beta}(\partial B),
\]
is defined on $\mathbb R^3$ and vanishes near the origin.  Set
\[
 \Phi_\rho=\operatorname{Id}+E\rho,
 \qquad \Omega_\rho=\Phi_\rho(B).
\]
For $\rho$ in a sufficiently small open ball $\mathcal U$ about zero in
$C^{2,\beta}(\partial B)$, $\Phi_\rho$ is a global $C^{2,\beta}$
diffeomorphism, equal to the identity outside a fixed ball, and
\[
 \partial\Omega_\rho
   =\{(1+\rho(\omega))\omega:\omega\in\partial B\}.
\]
We choose $\mathcal U$ convex and small enough that, uniformly for
$\rho\in\mathcal U$,
\begin{equation}
 \frac12|x-y|\leq|\Phi_\rho(x)-\Phi_\rho(y)|
                  \leq\frac32|x-y|,
 \qquad x,y\in\overline B.
 \label{eq:state-bilip}
\end{equation}
Introduce the fixed Banach spaces
\[
 \mathcal X_\beta=\{v\in C^{2,\beta}(\overline B):v|_{\partial B}=0\},
 \qquad \mathcal Y_\beta=C^{0,\beta}(\overline B).
\]
Here $\mathcal X_\beta$ is the closed zero-trace subspace; no closure of compactly
supported functions in the $C^{2,\beta}$ norm is intended.

Write $F_\rho=D\Phi_\rho$, $J_\rho=\det F_\rho$, and define
\begin{align}
 \mathcal L_\rho v
   &=-J_\rho^{-1}\operatorname{div}
       (J_\rho F_\rho^{-1}F_\rho^{-T}\nabla v),
       \label{eq:state-pullback-laplace}\\
 \mathcal C_\rho f(x)
   &=\int_B\frac{f(y)J_\rho(y)}
                      {|\Phi_\rho(x)-\Phi_\rho(y)|}\,dy.
       \label{eq:state-pullback-coulomb}
\end{align}
Thus $\mathcal L_\rho(u\circ\Phi_\rho)=(-\Delta u)\circ\Phi_\rho$.

\begin{lemma}[Material operators]
\label{lem:material}
After reducing $\mathcal U$ if necessary, the maps
\[
 \rho\longmapsto\mathcal L_\rho
       \in\mathcal L(\mathcal X_\beta,\mathcal Y_\beta),
 \qquad
 \rho\longmapsto\mathcal C_\rho
       \in\mathcal L(L^\infty(B),W^{1,\infty}(B))
\]
are continuously Fr\'echet differentiable.  In particular the second map
is continuously differentiable with values in
$\mathcal L(L^\infty(B),\mathcal Y_\beta)$.
On a smaller ball, there is a constant $C$ such that
\begin{align}
 \|\mathcal C_{\rho+h}-\mathcal C_\rho
       -D\mathcal C_\rho[h]\|_{\mathcal L(L^\infty,W^{1,\infty})}
       &\leq C\|h\|_{C^{2,\beta}}^2,
       \label{eq:state-coulomb-remainder}\\
 \|D\mathcal C_{\rho+h}-D\mathcal C_\rho\|
 _{\mathcal L(C^{2,\beta},\mathcal L(L^\infty,W^{1,\infty}))}
       &\leq C\|h\|_{C^{2,\beta}},
       \label{eq:state-coulomb-continuity}
\end{align}
whenever the segment joining $\rho$ to $\rho+h$ lies in that ball.
Consequently
\[
 (\rho,v)\longmapsto\mathcal C_\rho(v^2)v
 \quad\hbox{is }C^1\hbox{ from }\mathcal U\times \mathcal X_\beta
                          \hbox{ to }\mathcal Y_\beta.
\]
\end{lemma}

\begin{proof}
Put $C_\rho=F_\rho^{-1}$.  Expanding the chain rule gives
\[
 \mathcal L_\rho v
  =-a_{kj}(\rho)\partial_k\partial_jv
       -b_j(\rho)\partial_jv,
 \quad
 a_{kj}=\sum_i(C_\rho)_{ki}(C_\rho)_{ji},
 \quad
 b_j=\sum_{k,i}(C_\rho)_{ki}\partial_k(C_\rho)_{ji}.
\]
Matrix inversion is smooth on the open set of invertible matrices in
the Banach algebra $C^{1,\beta}(\overline B)$.  The coefficients $a_{kj}$
therefore depend smoothly on $\rho$ in $C^{1,\beta}$, and $b_j$ in
$C^{0,\beta}$.  The displayed differential operator gives the first
assertion.  In particular, its dependence is in the operator norm
$\mathcal L(\mathcal X_\beta,\mathcal Y_\beta)$, with no variable composition operator
left to differentiate.

We give the singular-kernel estimates for the second assertion.  For
$x\ne y$, put $r=|x-y|$, $d=\Phi_\rho(x)-\Phi_\rho(y)$, and
$k_\rho(x,y)=|d|^{-1}$.  For directions $h_1,h_2$, let
$\eta_i=Eh_i$ and $a_i=\eta_i(x)-\eta_i(y)$.  Direct differentiation
with respect to $\rho$ gives
\begin{align*}
 Dk_\rho[h_1]&=-\frac{d\cdot a_1}{|d|^3},\\
 D^2k_\rho[h_1,h_2]
   &=3\frac{(d\cdot a_1)(d\cdot a_2)}{|d|^5}
        -\frac{a_1\cdot a_2}{|d|^3}.
\end{align*}
The bound $|a_i|\leq\|\eta_i\|_{C^1}r$ and
\eqref{eq:state-bilip} imply, for $m=0,1,2$,
\begin{equation}
 |D^mk_\rho[h_1,\ldots,h_m]|
      \leq Cr^{-1}\prod_{i=1}^m\|h_i\|_{C^{2,\beta}},
 \qquad
 |\nabla_xD^mk_\rho[h_1,\ldots,h_m]|
      \leq Cr^{-2}\prod_{i=1}^m\|h_i\|_{C^{2,\beta}}.
 \label{eq:state-kernel-bounds}
\end{equation}
For the second bound, differentiating a homogeneous kernel in $d$ costs
one power of $r$, while differentiating an increment $a_i$ replaces its
$O(r)$ bound by $\|D\eta_i\|_\infty$.  This proves the bound for each
term in the two formulas, as well as for $m=0$.  Empty products equal
one.

The Jacobian and its first two parameter derivatives are uniformly
bounded in $L^\infty(B)$ on the chosen ball.  Moreover
\[
 DJ_\rho[h]=J_\rho\operatorname{tr}(F_\rho^{-1}DEh).
\]
The estimates \eqref{eq:state-kernel-bounds} thus hold, with the same
singularity orders, for the kernel
$K_\rho(x,y)=J_\rho(y)k_\rho(x,y)$ and its first two parameter
derivatives.  In dimension three,
\begin{equation}
 \sup_{x\in B}\int_B(|x-y|^{-1}+|x-y|^{-2})\,dy<\infty.
 \label{eq:state-integrable-kernels}
\end{equation}
Consequently integration of these kernels against $f\in L^\infty(B)$
defines bounded maps into $W^{1,\infty}(B)$.  To justify the weak
spatial derivative, one may delete $|x-y|\leq\varepsilon$, integrate
by parts against a compactly supported test function in $x$, and let
$\varepsilon\downarrow0$.  The extra boundary integral is $O(\varepsilon)$,
because the kernel is $O(\varepsilon^{-1})$ on a sphere of area
$O(\varepsilon^2)$; the remaining integrals converge by
\eqref{eq:state-integrable-kernels}.  This applies also to the first
two parameter derivatives.

In particular, the candidate first derivative is
\begin{equation}
 D\mathcal C_\rho[h]f(x)
  =\int_B f(y)\left\{
       \frac{DJ_\rho[h](y)}{|d|}
       -\frac{J_\rho(y)d\cdot(Eh(x)-Eh(y))}{|d|^3}
                    \right\}\,dy.
 \label{eq:state-coulomb-derivative}
\end{equation}
Taylor's formula for $K_{\rho+th}(x,y)$, and also for its weak
$x$-gradient, has a second derivative bounded respectively by
$C\|h\|_{C^{2,\beta}}^2r^{-1}$ and
$C\|h\|_{C^{2,\beta}}^2r^{-2}$.  Integrating the remainder proves
\eqref{eq:state-coulomb-remainder}.  Applying the fundamental theorem
of calculus to $DK_{\rho+th}[\ell]$, using the same second derivative
bounds, gives
\[
 \|(D\mathcal C_{\rho+h}-D\mathcal C_\rho)[\ell]f\|_{W^{1,\infty}}
 \leq C\|h\|_{C^{2,\beta}}\|\ell\|_{C^{2,\beta}}\|f\|_\infty,
\]
which proves \eqref{eq:state-coulomb-continuity}.

The convexity and boundedness of $B$ give the continuous inclusion
$W^{1,\infty}(B)\hookrightarrow C^{0,\beta}(\overline B)$.
Finally $v\mapsto v^2$ is smooth into $L^\infty(B)$, and multiplication
is continuous on $C^{0,\beta}(\overline B)$.  In particular the state
derivative of the nonlocal term is
\[
 D_v\big(\mathcal C_\rho(v^2)v\big)[w]
       =\mathcal C_\rho(v^2)w+2v\mathcal C_\rho(vw).
\]
This proves the last assertion.
\end{proof}

\begin{lemma}[The normalised material state]
\label{lem:state}
There are a neighbourhood $\mathcal U_0\subset\mathcal U$ of zero and
$q_*>0$ with the following properties.  For
$\rho\in\mathcal U_0$ and $0\leq q<q_*$ there is a normalised positive
solution $u_{\rho,q}\in C^{2,\beta}(\overline\Omega_\rho)$ of
\begin{equation}
 -\Delta u_{\rho,q}+q\big(|\cdot|^{-1}*u_{\rho,q}^2\big)u_{\rho,q}
       =\lambda_{\rho,q}u_{\rho,q},\qquad
 u_{\rho,q}|_{\partial\Omega_\rho}=0,
 \qquad \int_{\Omega_\rho}u_{\rho,q}^2=1.
 \label{eq:state-hartree-equation}
\end{equation}
The map
\[
 (\rho,q)\longmapsto
 (v_{\rho,q},\lambda_{\rho,q})
 :=(u_{\rho,q}\circ\Phi_\rho,\lambda_{\rho,q})
 \quad\hbox{takes values in }\mathcal X_\beta\times\mathbb R
\]
and has a $C^1$ extension to an open neighbourhood of $(0,0)$.
It is the unique solution branch close to the normalised first
Dirichlet eigenpair of $B$.  The physical normal derivative, pulled
back to the sphere,
\begin{equation}
 g(\rho,q)=
 (\partial_{\nu_\rho}u_{\rho,q})\circ\Phi_\rho,
 \label{eq:state-flux-definition}
\end{equation}
is $C^1$ with values in $C^{1,\beta}(\partial B)$ and is strictly
negative on $\partial B$.
\end{lemma}

\begin{proof}
Let $\phi$ be the positive normalised first Dirichlet eigenfunction of
$B$, and let $\lambda_0=\pi^2$.  Consider
\[
 \mathcal F(\rho,q,v,\lambda)=
 \left(\mathcal L_\rho v+q\mathcal C_\rho(v^2)v-\lambda v,
                     \int_Bv^2J_\rho-1\right)
       \in \mathcal Y_\beta\times\mathbb R.
\]
Lemma~\ref{lem:material} shows that $\mathcal F$ is $C^1$ on an open
set with $q\in\mathbb R$.  At $(0,0,\phi,\lambda_0)$ its derivative
in the last two variables is
\begin{equation}
 (w,\mu)\longmapsto
 \left(({-\Delta-\lambda_0})w-\mu\phi,
                         2\int_B\phi w\right).
 \label{eq:state-augmented-inverse}
\end{equation}
This is a bounded isomorphism from $\mathcal X_\beta\times\mathbb R$ onto
$\mathcal Y_\beta\times\mathbb R$.  Indeed, for a prescribed pair $(f,s)$,
set
\[
 \mu=-\int_B f\phi,\qquad w=\frac{s}{2}\phi+w_\perp,
\]
where
\[
 (-\Delta-\lambda_0)w_\perp
       =f-\left(\int_Bf\phi\right)\phi,
 \qquad w_\perp|_{\partial B}=0,
 \qquad\int_B\phi w_\perp=0.
\]
The first Dirichlet eigenvalue is simple.  The Fredholm alternative
therefore gives the unique orthogonal solution, and global Dirichlet
Schauder estimates give a bounded inverse in the displayed H\"older
spaces; see \cite[Chapter~6 and Theorem~8.38]{GT2001}.
The Banach implicit-function theorem now gives a $C^1$ branch
$(v_{\rho,q},\lambda_{\rho,q})$.  The auxiliary negative values of $q$
serve only to obtain an open parameter neighbourhood; subsequent
variational assertions concern $q\geq0$.

The branch is positive after reducing the neighbourhood.  In fact
$\phi(x)\geq c\operatorname{dist}(x,\partial B)$ for some $c>0$.
This follows, for example, from
$\phi(r)=(2\pi)^{-1/2}\sin(\pi r)/r$ with its continuous value at
$r=0$.  Since $v-\phi$ vanishes on the sphere, integration along the
segment from $x$ to its nearest boundary point gives
\[
 |v(x)-\phi(x)|\leq\|\nabla(v-\phi)\|_\infty
                         \operatorname{dist}(x,\partial B).
\]
The $C^{2,\beta}$ continuity of the branch therefore gives
$v_{\rho,q}\geq(c/2)\operatorname{dist}(\cdot,\partial B)>0$ in $B$.
Likewise $\partial_n v_{\rho,q}<0$ on the sphere, since
$\partial_n\phi=-\sqrt{\pi/2}$ there and the convergence is in $C^1$.

For the flux, the exact transformation law for the unit normal is
\[
 \nu_\rho\circ\Phi_\rho
       =\frac{F_\rho^{-T}n}{|F_\rho^{-T}n|},\qquad n(x)=x
       \quad(x\in\partial B).
\]
The tangential gradient of the zero boundary trace of $v_{\rho,q}$
vanishes, so
\begin{equation}
 g(\rho,q)=|F_\rho^{-T}n|\,\partial_n v_{\rho,q}.
 \label{eq:state-material-flux}
\end{equation}
The first factor depends smoothly on $\rho$ in $C^{1,\beta}(\partial B)$,
and the normal trace is a bounded map
$\mathcal X_\beta\to C^{1,\beta}(\partial B)$.  This proves the stated $C^1$
regularity and negativity.
\end{proof}

All parameter derivatives above are taken after the fixed pullback.
In particular, \eqref{eq:state-material-flux} does not require
differentiating a composition operator on a H\"older diffeomorphism
group.  The fixed-domain construction is consistent with the
eigenfunction and boundary-map construction in
\cite[Propositions~3.2 and 5.1]{PacardSicbaldi2009}.

Uniqueness of the nonnegative minimising state on a connected domain
is established in \cite[Proposition~2.9]{MMR2025}.  The next argument
also identifies the positive critical branch just constructed with
that minimising state.

\begin{lemma}[Selection of the ground state]
\label{lem:selection}
For $\rho\in\mathcal U_0$ and $0\leq q<q_*$, the solution from
Lemma~\ref{lem:state} is the unique nonnegative normalised minimiser
defining $E_q(\Omega_\rho)$.  Every real normalised minimiser is either
$u_{\rho,q}$ or $-u_{\rho,q}$.
\end{lemma}

\begin{proof}
We first record that all Coulomb expressions used below are finite for
arbitrary competitors in $H^1_0(\Omega_\rho)$.  If $z$ is such a
competitor, its zero extension belongs to $H^1(\mathbb R^3)$ and hence
to $L^2(\mathbb R^3)\cap L^6(\mathbb R^3)$.  Interpolation gives
$z\in L^{12/5}$, so $z^2\in L^{6/5}$.  The
Hardy--Littlewood--Sobolev inequality gives
\begin{equation}
 |D(f_1,f_2)|\leq C\|f_1\|_{L^{6/5}}\|f_2\|_{L^{6/5}},
 \qquad f_1,f_2\in L^{6/5}(\mathbb R^3);
 \label{eq:state-hls}
\end{equation}
see \cite[Theorem~4.3]{LiebLoss2001}.  In particular the integral is
absolutely convergent, by applying the same inequality to $|f_1|$ and
$|f_2|$.

For real $f\in L^{6/5}(\mathbb R^3)$ one also has $D(f,f)\geq0$.
For $f\in C_c^\infty$, put $U=|\cdot|^{-1}*f$.  The identity
$-\Delta U=4\pi f$ and integration by parts on expanding balls give
\[
 D(f,f)=\frac1{4\pi}\int_{\mathbb R^3}|\nabla U|^2\geq0;
\]
the boundary term tends to zero because $U=O(|x|^{-1})$ and
$\nabla U=O(|x|^{-2})$ at infinity.  Approximation in $L^{6/5}$ and
\eqref{eq:state-hls} extend the nonnegativity to the stated class.
This is the Coulomb positivity also recorded in
\cite[Section~2, p.~1034]{MMR2025}.

Abbreviate $u=u_{\rho,q}$, $\lambda=\lambda_{\rho,q}$, and
$V_u=|\cdot|^{-1}*u^2$.  Lemma~\ref{lem:material} implies that $V_u$
is bounded on $\Omega_\rho$.  On this fixed connected domain, $u$ is
a positive eigenfunction of the self-adjoint Dirichlet operator
$-\Delta+qV_u$.  The principal Dirichlet eigenvalue is simple and
has a positive eigenfunction \cite[Theorem~8.38]{GT2001}.  Positivity
and orthogonality of eigenfunctions belonging to different eigenvalues
show that $\lambda$ is this principal eigenvalue.  Thus, for every
real $z\in H^1_0(\Omega_\rho)$ with $\|z\|_2=1$,
\[
 Q_u(z):=\int_{\Omega_\rho}|\nabla z|^2+qD(u^2,z^2)
         \geq\lambda=Q_u(u).
\]
Putting $\delta=z^2-u^2\in L^{6/5}$ and expanding the quadratic
Coulomb form gives the exact identity
\begin{equation}
 \mathcal E_q(z)-\mathcal E_q(u)
   =Q_u(z)-Q_u(u)+\frac q2D(\delta,\delta)\geq0.
 \label{eq:state-ground-state-identity}
\end{equation}
Here the final inequality uses precisely $q\geq0$.  It proves that
the branch is a global minimiser over all normalised states on the
fixed domain.  Equality forces equality in the principal Rayleigh
inequality, and simplicity gives $z=\pm u$.  This proves both claims.
\end{proof}

\begin{lemma}[First variation and constant boundary flux]
\label{lem:variation}
Fix $0\leq q<q_*$.  If $t\mapsto\rho_t\in\mathcal U_0$ is a $C^1$
curve, set $\rho=\rho_0$, $\Omega=\Omega_\rho$, $u=u_{\rho,q}$, and
\[
 T_t=\Phi_{\rho_t}\circ\Phi_\rho^{-1},\qquad
 V=\left.\frac{d}{dt}\right|_{t=0}T_t.
\]
Then
\begin{equation}
 \left.\frac{d}{dt}\right|_{t=0}E_q(\Omega_{\rho_t})
       =-\int_{\partial\Omega}(\partial_\nu u)^2V\cdot\nu\,dS.
 \label{eq:state-hadamard}
\end{equation}
If $\Omega_\rho$ is stationary for $E_q$ under the volume constraint
within this chart, then $\partial_\nu u$ is a negative constant on
$\partial\Omega_\rho$.  In particular this conclusion holds for a
volume-constrained minimising domain lying in the chart.
\end{lemma}

\begin{proof}
By Lemmas~\ref{lem:state} and \ref{lem:selection}, the material curve
$u_{\rho_t,q}\circ T_t$ is $C^1$ in
$C^{2,\beta}(\overline\Omega)$ and consists of the normalised
minimisers.  Denote its derivative by $m$ and put
\[
 w=m-\nabla u\cdot V\in C^{1,\beta}(\overline\Omega).
\]
The material trace is zero for every $t$.  Consequently
\begin{equation}
 w=-\partial_\nu u\,V\cdot\nu\quad\hbox{on }\partial\Omega,
 \qquad \int_\Omega uw=0.
 \label{eq:state-shape-normalisation}
\end{equation}
The second identity follows by differentiating normalisation:
\[
 0=\int_\Omega(2um+u^2\operatorname{div}V)
   =2\int_\Omega uw+\int_{\partial\Omega}u^2V\cdot\nu
   =2\int_\Omega uw.
\]

For clarity, the transport differentiation of the Coulomb term can
be justified without a boundary distribution.  Let $f_t$ denote
the zero extension of $u_{\rho_t,q}^2$.  Then
\begin{equation}
 \frac{f_t-f_0}{t}\longrightarrow 2uw\,\mathbf1_\Omega
       \quad\hbox{in }L^{6/5}(\mathbb R^3).
 \label{eq:state-density-derivative}
\end{equation}
Here is a direct verification.  Extend the material states from the
fixed $C^{2,\beta}$ domain by a fixed bounded extension operator and
use $T_t$ in a fixed neighbourhood.  Taylor's formula, first in the
material parameter and then in the spatial variable, gives the
Eulerian expansion $u_t=u+tw+o(t)$ uniformly on the common domain.
All supports lie in one bounded set.  The symmetric difference of
the two domains is contained in a boundary strip of width $O(|t|)$
and measure $O(|t|)$.  On that strip, the state belonging to its
domain is $O(|t|)$, because the states have a uniform Lipschitz bound
and vanish at their respective boundaries.  Its squared density
divided by $t$ is therefore $O(|t|)$.  The same is true of $uw$ on
the part of the strip in $\Omega$.  These estimates prove
\eqref{eq:state-density-derivative}, together with the uniform
interior expansion.  By the bilinear continuity
\eqref{eq:state-hls},
\begin{equation}
 \left.\frac{d}{dt}\right|_{t=0}\frac q2D(f_t,f_t)
                =2qD(u^2,uw)=2q\int_\Omega V_uuw.
 \label{eq:state-hartree-variation}
\end{equation}

The material formula for the Dirichlet term is
\[
 \left.\frac{d}{dt}\right|_{t=0}\int_{\Omega_{\rho_t}}
                        |\nabla u_{\rho_t,q}|^2
 =\int_\Omega\bigl(2\nabla u\cdot\nabla m
          -2\nabla u\cdot(DV)^T\nabla u
          +|\nabla u|^2\operatorname{div}V\bigr).
\]
Using
$\nabla m-(DV)^T\nabla u=\nabla w+D^2u\,V$ and the divergence
theorem, this becomes
\[
 2\int_\Omega\nabla u\cdot\nabla w
           +\int_{\partial\Omega}|\nabla u|^2V\cdot\nu.
\]
Adding \eqref{eq:state-hartree-variation}, integrating by parts, and
then using \eqref{eq:state-hartree-equation} and
\eqref{eq:state-shape-normalisation}, we obtain
\begin{align*}
 \left.\frac{d}{dt}\right|_0E_q(\Omega_{\rho_t})
 &=2\int_\Omega(-\Delta u+qV_uu)w
     +2\int_{\partial\Omega}\partial_\nu u\,w
     +\int_{\partial\Omega}|\nabla u|^2V\cdot\nu\\
 &=-\int_{\partial\Omega}(\partial_\nu u)^2V\cdot\nu.
\end{align*}
In the last step, the tangential gradient of $u$ is zero on the
boundary.  This proves \eqref{eq:state-hadamard}.

It remains to extract the volume-constrained condition without
requiring extra regularity of the graph parametrisation.  For the
radial graph,
\[
 |\Omega_\rho|=\frac13\int_{\partial B}(1+\rho)^3\,dS,
 \qquad
 D|\Omega_\rho|[h]=\int_{\partial B}(1+\rho)^2h\,dS.
\]
The radial normal velocity and surface Jacobian satisfy
$ (V\cdot\nu)\,dS_{\partial\Omega}
      =h(1+\rho)^2\,dS_{\partial B}$.
Thus, with $g=g(\rho,q)$,
\[
 DE_q(\Omega_\rho)[h]
       =-\int_{\partial B}g^2(1+\rho)^2h\,dS.
\]
Every $h\in C^{2,\beta}(\partial B)$ with
$\int(1+\rho)^2h=0$ is tangent to an exactly volume-preserving
curve: use $\rho_t=\rho+th+a(t)$ and solve the scalar volume equation
for $a(t)$ by the implicit-function theorem.  Its derivative in
$a$ is $\int(1+\rho)^2>0$, and $a'(0)=0$.
Stationarity therefore implies that the linear functional with
density $g^2(1+\rho)^2$ is proportional to that with density
$(1+\rho)^2$.  Testing against smooth functions and using
$(1+\rho)^2>0$ shows that $g^2$ is constant.  Finally
$g<0$ by Lemma~\ref{lem:state}, so $g$ itself is a negative constant.
\end{proof}

\section{Local rigidity of the ball}\label{sec:local}

Fix $0<\beta<1$. Write $\Pi$ for the $L^2(\mathbb S^2)$ projection that removes the constant functions and the three coordinate functions. The same finite-rank formula defines a bounded projection on each H\"older space. Set
\[
 X_\beta=\Pi C^{2,\beta}(\mathbb S^2),\qquad
 Y_\beta=\Pi C^{1,\beta}(\mathbb S^2).
\]
Orthogonality throughout this section uses the standard surface measure. We first choose a genuine volume and translation slice, rather than merely imposing linearised constraints.

\begin{lemma}[Volume and barycentre slice]\label{lem:gauge}
For $h$ in a neighbourhood of zero in $X_\beta$, there are unique smooth functions $c(h)\in\mathbb R$ and $b(h)\in\mathbb R^3$, close to zero, such that
\[
 \rho(h)(\omega)=h(\omega)+c(h)+b(h)\cdot\omega
\]
defines a radial domain $\Omega_{\rho(h)}$ of volume $V=4\pi/3$ and barycentre zero. They satisfy $c(0)=b(0)=0$ and $Dc(0)=Db(0)=0$. Every sufficiently small radial graph with these two geometric constraints has this form, with $h=\Pi\rho$.
\end{lemma}
\begin{proof}
For a radial graph the volume and first moment are
\[
 \frac13\int_{\mathbb S^2}(1+\rho)^3\,d\sigma,
 \qquad \frac14\int_{\mathbb S^2}(1+\rho)^4\omega\,d\sigma.
\]
These are smooth polynomial maps. At $\rho=0$, their derivative on $c+b\cdot\omega$ is $(4\pi c,(4\pi/3)b)$. Their derivative on $h\in X_\beta$ is zero. The finite-dimensional implicit-function theorem, with Banach parameter $h$, proves every assertion, including the converse by uniqueness.
\end{proof}

Let $g(\rho,q)$ be the physical outward normal derivative of the positive normalised state, pulled back to $\mathbb S^2$ along $\omega\mapsto(1+\rho(\omega))\omega$. Section~\ref{sec:state} proves the $C^1$ map
\[
 g:U\times(-q_*,q_*)\subset C^{2,\beta}(\mathbb S^2)\times\mathbb R
       \longrightarrow C^{1,\beta}(\mathbb S^2).
\]
For $q\geq0$ this state is the ground state selected by the energy infimum. Negative $q$ is used only to have an open parameter neighbourhood for the implicit-function theorem. Define
\[
 F(q,h)=\Pi g(\rho(h),q)\in Y_\beta.
\]
On the ball, rotation invariance and local uniqueness of the state branch imply that the state is radial, also for the auxiliary signed parameters. Thus $F(q,0)=0$. A domain whose physical normal derivative is constant necessarily satisfies $F(q,h)=0$. We do not need the converse implication.

\subsection{The derivative at zero charge}
The first normalised Dirichlet eigenfunction and its eigenvalue on $B$ are
\[
 \phi(r)=\frac{1}{\sqrt{2\pi}}\frac{\sin(\pi r)}{r},
 \qquad \lambda_0=\pi^2.
\]
The value at the origin is interpreted by continuity. Put $a=\pi/\sqrt{2\pi}$; then $\phi'(1)=-a$ and $\phi''(1)=2a$. For $h\in X_\beta$, the Eulerian derivative $w$ of this normalised eigenfunction under the boundary displacement $h\omega$ obeys
\begin{equation}\label{eq:eulerian}
 (-\Delta-\pi^2)w=\dot\lambda\phi\quad\hbox{in }B,
 \qquad w=ah\quad\hbox{on }\mathbb S^2,
 \qquad \int_B\phi w=0.
\end{equation}
These identities follow by differentiating the fixed-domain equation of Lemma~\ref{lem:state}, and then subtracting the material transport $\nabla\phi\cdot\dot\Phi$; at the smooth reference ball all quantities are classical. Green's identity with $\phi$ gives
\[
 \dot\lambda=\int_{\mathbb S^2}w\partial_\nu\phi\,d\sigma
             =-a^2\int_{\mathbb S^2}h\,d\sigma=0.
\]
The constraint $\int\phi w=0$ follows from normalisation because $\phi$ vanishes on the boundary. The derivatives of $c,b$ in Lemma~\ref{lem:gauge} vanish, so they add no first-order term here.

Differentiating the physical normal derivative gives
\begin{equation}\label{eq:flux-linearisation}
 D_hF(0,0)h=\Pi\bigl(\partial_r w+2ah\bigr).
\end{equation}
Indeed the derivative of the normal is tangential and hence orthogonal to the radial vector $\nabla\phi$; the displacement of the evaluation point contributes $\phi''(1)h$. This formula keeps the physical normal, rather than replacing it with a coordinate derivative away from the ball.

\begin{lemma}[All-mode nondegeneracy]\label{lem:modes}
Let $Y_l$ be any spherical harmonic of degree $l\geq1$. The unique solution of
\[
 (-\Delta-\pi^2)W=0,\quad W|_{\mathbb S^2}=Y_l,
 \quad \int_B\phi W=0
\]
is $R_l(r)Y_l(\omega)$ with $R_l(1)=1$. If $\kappa_l=R_l'(1)$, then $\kappa_1=-2$ and $\kappa_l>-2$ for every $l\geq2$. Consequently the derivative in \eqref{eq:flux-linearisation} has multiplier $a(\kappa_l+2)>0$ on every degree $l\geq2$.
\end{lemma}
\begin{proof}
For $l\geq1$ consider the radial energy
\[
 J_l(R)=\int_0^1\left(r^2(R')^2+l(l+1)R^2-\pi^2r^2R^2\right)\,dr.
\]
Use the common radial energy space defined by $\int_0^1(r^2(R')^2+R^2)\,dr<\infty$, with trace $R(1)=1$. Multiplication by a fixed degree-$l$ spherical harmonic identifies this with the corresponding $H^1(B)$ sector. On zero-boundary differences the quadratic form is coercive: the sector is orthogonal to the first eigenfunction, and the second Dirichlet eigenvalue of $B$ is strictly greater than $\pi^2$. The direct method therefore gives a unique minimiser $R_l$. The Euler--Lagrange equation and interior regularity show it is the regular radial solution
\[
 -(r^2R_l')'+l(l+1)R_l-\pi^2r^2R_l=0.
\]
The singular solution at the origin has infinite energy. Integration by parts, or the weak equation with an extension of the boundary value, yields $J_l(R_l)=R_l'(1)=\kappa_l$. The regular solution is also $j_l(\pi r)/j_l(\pi)$, but no estimate on a finite list of Bessel values is used. Coercivity excludes $j_l(\pi)=0$.

Differentiating the radial equation for $\phi$ shows $R_1=-\phi'/a$. Hence $R_1(1)=1$ and $\kappa_1=-\phi''(1)/a=-2$. For $l\geq2$, $R_l$ is an admissible competitor for $J_1$, so
\[
 \kappa_l=J_1(R_l)+(l(l+1)-2)\int_0^1R_l^2\,dr
 \geq\kappa_1+(l(l+1)-2)\int_0^1R_l^2\,dr>\kappa_1.
\]
The final integral is positive because the trace is one. Equation~\eqref{eq:flux-linearisation} now gives the stated multiplier. The vanishing degree-one multiplier is exactly the translation freedom removed by the gauge.
\end{proof}

\begin{lemma}[H\"older isomorphism]\label{lem:inverse}
The map $L=D_hF(0,0):X_\beta\to Y_\beta$ is a bounded isomorphism.
\end{lemma}
\begin{proof}
Let $\Lambda_0$ denote the harmonic Dirichlet-to-Neumann operator on $B$. The Robin problem
\[
 \Delta H=0\quad\hbox{in }B,\qquad
 \partial_\nu H+2H=f\quad\hbox{on }\mathbb S^2
\]
has a unique solution for $f\in C^{1,\beta}(\mathbb S^2)$, with boundary trace in $C^{2,\beta}$ and a bounded Schauder estimate. The applicable oblique boundary results are \cite[Theorems~6.30--6.31]{GT2001}: here the oblique vector is the outward normal, its normal component is one, and the boundary zeroth-order coefficient is $2>0$. Uniqueness also follows by integrating $|\nabla H|^2+2\int_{\mathbb S^2}H^2=\int_{\mathbb S^2}fH$. Since the operator commutes with the spherical harmonic projections, $\Lambda_0+2$ is an isomorphism from $X_\beta$ to $Y_\beta$.

Let $\Lambda_\pi$ be the reduced Helmholtz Dirichlet-to-Neumann operator on these spaces. It is well-defined: for boundary data $h$ with vanishing degree zero, the Fredholm compatibility condition at the eigenvalue $\pi^2$ holds, and $\int_B\phi W=0$ fixes the eigenfunction ambiguity. The absence of degree one is preserved as well. If $H$ is the harmonic extension of $h$ and $W$ its reduced Helmholtz extension, then $Z=W-H$ satisfies
\[
 (-\Delta-\pi^2)Z=\pi^2H,\qquad Z|_{\mathbb S^2}=0,
 \qquad \int_B\phi Z=0.
\]
All data have degrees at least two. The global Dirichlet estimate \cite[Theorem~6.6]{GT2001} and its higher-order boundary regularity version \cite[Theorem~6.19, with $k=2$]{GT2001}, applied on the smooth ball to the constant-coefficient operator, give the reduced-resolvent estimate
\[
 \|Z\|_{C^{4,\beta}(\overline B)}
       \leq C\|H\|_{C^{2,\beta}(\overline B)}
       \leq C'\|h\|_{C^{2,\beta}(\mathbb S^2)}.
\]
The usual lower-order term in the Schauder estimate is controlled by invertibility on the orthogonal complement of $\phi$. Thus $\Lambda_\pi-\Lambda_0$, the boundary trace of $\partial_\nu Z$, is bounded into $C^{3,\beta}$ and compact into $C^{1,\beta}$. The compactness follows, for example, from the compact inclusion $C^{3,\beta}(\mathbb S^2)\hookrightarrow C^{1,\beta}(\mathbb S^2)$.

Formula~\eqref{eq:flux-linearisation} says $L=a(\Lambda_\pi+2)$ on $X_\beta$. It is a compact perturbation of an isomorphism, hence Fredholm of index zero. If $Lh=0$, project onto each spherical harmonic. Lemma~\ref{lem:modes} forces all coefficients of $h$ to vanish. Completeness in $L^2(\mathbb S^2)$ gives $h=0$; no convergence claim for spectral partial sums in a H\"older norm is needed. The Fredholm alternative and bounded inverse theorem finish the proof.
\end{proof}

\begin{theorem}[Local constant-flux rigidity]\label{thm:local}
For every fixed $0<\beta<1$, there are $\eta>0$ and $q_{\mathrm{loc}}>0$ with the following property. Suppose $0\leq q<q_{\mathrm{loc}}$ and $\Omega_\rho$ is a radial domain with
\[
 \rho\in C^{2,\beta}(\mathbb S^2),\quad
 \|\rho\|_{C^{2,\beta}}<\eta,\quad
 |\Omega_\rho|=V,\quad \int_{\Omega_\rho}x\,dx=0.
\]
If the positive normalised ground state for $E_q(\Omega_\rho)$ has constant physical outward normal derivative, then $\rho=0$.
\end{theorem}
\begin{proof}
Lemmas~\ref{lem:material}--\ref{lem:selection} and \ref{lem:gauge} make $F$ a $C^1$ map on an open neighbourhood of $(0,0)$ in $\mathbb R\times X_\beta$. Lemma~\ref{lem:inverse} makes its $h$ derivative there invertible. The Banach implicit-function theorem supplies a neighbourhood in which the zeros of $F(q,h)$ form a unique branch. Since $F(q,0)=0$, that branch is identically zero. Choose $\eta,q_{\mathrm{loc}}$ so that the gauge of every domain in the statement lies in this neighbourhood and the ground-state selection lemma applies. Constant physical flux implies $F(q,h)=0$, so $h=0$, and the gauge uniqueness gives $\rho=0$.
\end{proof}

The spectral mechanism in this section is classical; compare \cite[Propositions 3.2, 4.2--4.3 and Section 5]{PacardSicbaldi2009}. The Hartree-specific input is the differentiable self-consistent state and physical-flux map of Section~\ref{sec:state}. Theorem~\ref{thm:local} by itself does not classify arbitrary competitors. That step uses the global localisation theorem and the representative arguments in the next section.

\section{Passage from local rigidity to the global problem}
\label{sec:global}

We use the spaces $X(\Omega)$ and $H^1_0(\Omega)$ and the energies
$E_q^{\mathrm{ae}}(\Omega)$ and $E_q(\Omega)$ defined in the introduction.
The former convention imposes vanishing almost everywhere outside a
measurable domain; the latter imposes vanishing quasi-everywhere outside
a quasi-open domain, as in \cite[Section~2.3, pp.~1037--1038]{MMR2025}.
Throughout this section, $B=B_1(0)$, $V=|B|$, and $S=\partial B$.

These conventions agree on a bounded Lipschitz domain $D$:
\begin{equation}
 X(D)=H^1_0(D).
 \label{eq:global-lipschitz-support}
\end{equation}
Indeed, the zero extension of an element of $H^1_0(D)$ belongs to
$X(D)$.  Conversely, a global $H^1$ function that vanishes a.e. outside
$D$ has zero exterior trace.  The interior and exterior traces of a
global $H^1$ function agree across a Lipschitz boundary, so its interior
trace is zero; the trace characterisation of $H^1_0(D)$ proves the
reverse inclusion.  Thus both its Sobolev and its quasi-continuous zero
extension vanish q.e. outside $D$.  No analogous assertion will be made
for an arbitrary quasi-open set that merely agrees a.e. with $D$.

\begin{lemma}[Finite-volume domains and positivity supports]
\label{lem:measurable}
Fix $q>0$.  If $0<|\Omega|<\infty$, the embedding
$X(\Omega)\hookrightarrow L^2(\mathbb R^3)$ is compact.  The same is
true of $H^1_0(\Omega)$ when $\Omega$ is quasi-open.  Each of the two
energy infima, when finite, is attained by a nonnegative normalised
state.

Moreover, the global infima
\begin{equation}
 \begin{split}
 I_q^{\mathrm{ae}}&=
   \inf_{\substack{\Omega\text{ measurable}\\|\Omega|=V}}
           E_q^{\mathrm{ae}}(\Omega),\\
 I_q&=\inf_{\substack{\Omega\text{ quasi-open}\\|\Omega|=V}}
           E_q(\Omega)
 \end{split}
 \label{eq:global-infima}
\end{equation}
agree.  If $\Omega$ attains either global infimum and $u\geq0$ is a
normalised minimising state, then
\[
 Q=\{\widetilde u>0\}
\]
is a quasi-open global minimiser for $I_q$, has measure $V$, and agrees
a.e. with $\Omega$.  In the quasi-open convention, $Q\subset\Omega$
q.e. as well.
\end{lemma}

\begin{proof}
For every $u\in X(\Omega)$, the Sobolev inequality gives
\begin{equation}
 \int_{\mathbb R^3\setminus B_R}|u|^2
 \leq |\Omega\setminus B_R|^{2/3}\|u\|_6^2
 \leq C|\Omega\setminus B_R|^{2/3}\|\nabla u\|_2^2.
 \label{eq:global-tail}
\end{equation}
The last coefficient tends to zero as $R\to\infty$, since $|\Omega|$
is finite.  A bounded sequence in $X(\Omega)$ has a subsequence
converging in $L^2(B_R)$ for every fixed $R$, by local Rellich
compactness and a diagonal extraction.  Estimate
\eqref{eq:global-tail} makes this subsequence Cauchy in global $L^2$.
This proves compactness without any boundedness assumption on
$\Omega$.  The space $X(\Omega)$ is a closed linear subspace of $H^1$.
The quasi-open space $H^1_0(\Omega)$ is also closed: a strongly
$H^1$-convergent sequence has a subsequence whose quasi-continuous
representatives converge q.e., which preserves vanishing q.e. on the
complement.  Closed linear subspaces are weakly closed, and the same
compactness proof applies to $H^1_0(\Omega)\subset X(\Omega)$.

For completeness, the Coulomb term is finite on all of $H^1(\mathbb
R^3)$.  The translated Hardy inequality and Cauchy--Schwarz imply,
for each $a\in\mathbb R^3$,
\[
 \int_{\mathbb R^3}\frac{|u(x)|^2}{|x-a|}\,dx
 \leq\|u\|_2\left\|\frac{u}{|\cdot-a|}\right\|_2
 \leq 2\|u\|_2\|\nabla u\|_2.
\]
Consequently
\begin{equation}
 0\leq D(u^2,u^2)\leq2\|u\|_2^3\|\nabla u\|_2<\infty.
 \label{eq:global-coulomb-bound}
\end{equation}
This is the estimate of \cite[Lemma~2.1, p.~1033]{MMR2025}, whose
zero-extension proof applies unchanged to global $H^1$ functions.

Take a normalised minimising sequence in either support space.
Nonnegativity of the two energy terms bounds it in $H^1$.
Weak closure and the compactness just proved give a weak $H^1$ and
strong $L^2$ limit $u$ in the same space, with $\|u\|_2=1$.
After a further subsequence, $u_n\to u$ a.e.; hence
$u_n(x)^2u_n(y)^2\to u(x)^2u(y)^2$ for a.e. $(x,y)$.
Fatou's lemma for the nonnegative Coulomb integrand, together with weak
lower semicontinuity of the Dirichlet integral, proves attainment.
Replacing the state by its absolute value preserves admissibility,
normalisation, and energy, so the minimiser can be chosen nonnegative.

For any nonnegative normalised admissible state $u$, its
quasi-continuous positivity set $Q=\{\widetilde u>0\}$ is quasi-open,
and $u\in H^1_0(Q)$.  Also $Q\subset\Omega$ a.e. and
$0<m:=|Q|\leq V$ when $|\Omega|=V$.  Put
$s=(V/m)^{1/3}\geq1$ and $u_s(x)=s^{-3/2}u(x/s)$.  Dilation
preserves quasi-openness and q.e. support, and
\begin{equation}
 |sQ|=V,\qquad \|u_s\|_2=1,\qquad
 \mathcal E_q(u_s)
 =s^{-2}\int|\nabla u|^2+\frac q2s^{-1}D(u^2,u^2)
 \leq\mathcal E_q(u).
 \label{eq:global-dilation}
\end{equation}
The inequality is strict if $s>1$: a nonzero normalised $H^1(\mathbb
R^3)$ function cannot have zero Dirichlet integral.  Since $u_s$ is an
admissible state on the quasi-open set $sQ$, formula
\eqref{eq:global-dilation} yields $I_q\leq\mathcal E_q(u)$ for every
state in the measurable formulation.  Taking the relevant infima
shows $I_q\leq I_q^{\mathrm{ae}}$.  The reverse inequality follows
from $H^1_0(\Omega)\subset X(\Omega)$ and inclusion of quasi-open
domains among measurable domains.

If $\Omega$ is a global minimiser, take an attaining state $u$ as above.
Strictness in \eqref{eq:global-dilation} rules out $m<V$, since it would
produce a competitor below the common global infimum.  Thus $m=V$.
It follows that $\Omega=Q$ a.e. and
\[
 I_q\leq E_q(Q)\leq\mathcal E_q(u)=I_q.
\]
Finally, if $u\in H^1_0(\Omega)$ in the quasi-open convention, its
quasi-continuous representative is zero q.e. off $\Omega$, giving the
stated q.e. inclusion.
\end{proof}

We next record precisely the published geometric input.
Theorem~1.2 of \cite[p.~1031]{MMR2025} states that, for every
$\varepsilon>0$, there exist $q_*(\varepsilon)>0$ and
$\alpha\in(0,1)$ such that, for every $0<q\leq q_*(\varepsilon)$,
the volume-constrained problem has an optimal set and every optimal
set is $C^{2,\alpha}$-nearly spherical with graph norm at most
$\varepsilon$.  We use this theorem in its quasi-open formulation,
with the regular representatives supplied by its proof.  Explicitly,
a minimising positivity domain has, after translation, a bounded open
representative $D$ with
\begin{equation}
 \partial D=\{(1+\varphi(\theta))\theta:\theta\in S\},
 \qquad
 \|\varphi\|_{C^{2,\alpha}(S)}\leq\varepsilon,
 \qquad |D|=V.
 \label{eq:global-source-graph}
\end{equation}
The existence and regularity in this input are entirely those of
\cite{MMR2025}, not new assertions here.  The positivity formulation is
explained in its Section~2.3, Lemmas~3.6 and~3.11, and Theorem~4.4;
Section~5 removes the equiboundedness restriction.  It is enough below
to use that the representative agrees a.e. with the original domain.
Indeed, its associated minimising state then belongs to $X(D)$, hence
to $H^1_0(D)$ by \eqref{eq:global-lipschitz-support}, so $D$ is itself
a global minimiser.  Equality up to capacity will be established
separately at the end of the proof.

The order of quantifiers in this input permits $\alpha$ to depend on
$\varepsilon$; the bounded-domain statement
\cite[Theorem~3.1, p.~1038]{MMR2025} explicitly writes
$\alpha=\alpha(\varepsilon)$.  The following lemma supplies the fixed
H\"older topology needed in Theorem~\ref{thm:local}.

\begin{lemma}[Entrance into one fixed H\"older neighbourhood]
\label{lem:entrance}
There exist $\beta\in(0,1)$ and $\bar q>0$ with the following property.
For every sufficiently small $q\in(0,\bar q]$, every regular global
minimiser, translated to have barycentre zero, can be written uniquely
as
\[
 D_\rho=\{r\theta:\theta\in S,\ 0\leq r<1+\rho(\theta)\},
 \qquad \rho\in C^{2,\beta}(S).
\]
For every $\eta>0$ there is $q_\eta\in(0,\bar q]$ such that all these
graphs satisfy $\|\rho\|_{C^{2,\beta}(S)}<\eta$ whenever
$0<q<q_\eta$.
\end{lemma}

\begin{proof}
Fix a small number $\bar\varepsilon>0$ once and for all.
Applying \eqref{eq:global-source-graph} at this one tolerance gives
$\bar q>0$ and $\bar\alpha\in(0,1)$ such that every regular minimiser
for $0<q\leq\bar q$ has some translated graph $D_\varphi$ with
\begin{equation}
 \|\varphi\|_{C^{2,\bar\alpha}(S)}\leq\bar\varepsilon.
 \label{eq:global-fixed-bound}
\end{equation}
Choose $\bar\varepsilon$ small enough for the following elementary
recentring estimates.  We supply them because the centres selected
by the source theorem at different tolerances need not agree.

Let $b$ be the barycentre of $D_\varphi$.  Polar integration gives
\begin{equation}
 b=\frac1{4V}\int_S(1+\varphi(\theta))^4\theta\,d\sigma(\theta),
 \qquad |b|\leq C\|\varphi\|_{C^0(S)},
 \label{eq:global-barycentre}
\end{equation}
since $\int_S\theta\,d\sigma=0$.
Define
\[
 T(\theta)=\frac{(1+\varphi(\theta))\theta-b}
                 {|(1+\varphi(\theta))\theta-b|}.
\]
Its denominator is bounded away from zero, and its $C^1$ distance
from the identity is bounded by $C\|\varphi\|_{C^1}$.
For $\bar\varepsilon$ sufficiently small, $T$ is a local
diffeomorphism homotopic to the identity, hence a degree-one covering
and a global diffeomorphism of $S$.  In the new angular variable
$\omega=T(\theta)$ the centred radial function is
\begin{equation}
 1+\rho(\omega)=
 \left|(1+\varphi(T^{-1}(\omega)))T^{-1}(\omega)-b\right|.
 \label{eq:global-recentre}
\end{equation}
In particular, $D_\varphi-b=D_\rho$.

The ordinary chain rule in a fixed finite atlas of $S$, together with
\eqref{eq:global-barycentre} and the inverse-function derivative
formulas, gives
\begin{equation}
 \|\rho\|_{C^2}\leq C\|\varphi\|_{C^2},
 \qquad
 \|\rho\|_{C^{2,\bar\alpha}}\leq
                  C_{\bar\alpha}\|\varphi\|_{C^{2,\bar\alpha}}
 \label{eq:global-recentre-estimates}
\end{equation}
for graphs satisfying \eqref{eq:global-fixed-bound}.
Here the first estimate holds throughout a fixed small $C^2$
neighbourhood and does not involve a H\"older exponent.
To see the uniformity of the second estimate, $DT$ and its inverse
are uniformly bounded, while
$D(T^{-1})=(DT)^{-1}\circ T^{-1}$ and the formula for
$D^2(T^{-1})$ involves only $D^2T$ and three bounded inverse
first-derivative factors.  H\"older product and composition estimates
then bound these quantities by the fixed $C^{2,\bar\alpha}$ bound.
In \eqref{eq:global-recentre}, after subtracting the constant radius
one, each $C^2$ term contains $b$, $\varphi$, or a derivative of
$\varphi$, which proves the first estimate and its smallness.
These are estimates for individual coordinate changes; no
Fr\'echet differentiability of composition or inversion on a H\"older
diffeomorphism space is asserted or needed.

It follows from \eqref{eq:global-fixed-bound} and
\eqref{eq:global-recentre-estimates} that there is a fixed $M<\infty$
with
\begin{equation}
 \|\rho\|_{C^{2,\bar\alpha}}\leq M
 \quad\text{for every minimiser and }0<q\leq\bar q.
 \label{eq:global-centred-bound}
\end{equation}
Now apply the source theorem a second time, at an arbitrary smaller
tolerance $\varepsilon$.  Its possibly different exponent is
irrelevant to the consequence $\|\varphi\|_{C^2}\leq\varepsilon$.
By the first estimate in \eqref{eq:global-recentre-estimates}, the
barycentrically centred graph has $C^2$ norm at most
$C\varepsilon$.  This is the same graph $\rho$ as in
\eqref{eq:global-centred-bound}: the barycentre is intrinsic, and a
star-shaped domain has only one radial function about a given centre.
Therefore
\begin{equation}
 \sup_{\substack{0<q\leq\delta\\D_\rho\text{ centred global minimiser}}}
       \|\rho\|_{C^2(S)}\longrightarrow0
 \qquad(\delta\downarrow0).
 \label{eq:global-C2-small}
\end{equation}
Set, for example, $\beta=\bar\alpha/2$.  H\"older interpolation on
$S$ gives
\[
 \|\rho\|_{C^{2,\beta}}
 \leq C\left(\|\rho\|_{C^2}
       +\|\rho\|_{C^2}^{1-\beta/\bar\alpha}
        \|\rho\|_{C^{2,\bar\alpha}}^{\beta/\bar\alpha}\right).
\]
This follows by applying
$[f]_\beta\leq C\|f\|_\infty^{1-\beta/\bar\alpha}
[f]_{\bar\alpha}^{\beta/\bar\alpha}$ to second derivatives in a
fixed atlas.  Equations \eqref{eq:global-centred-bound} and
\eqref{eq:global-C2-small} prove the assertion uniformly over all
minimisers.  In particular, $\beta$ is chosen before the local
rigidity neighbourhood is fixed.
\end{proof}

Lemma~\ref{lem:entrance} places every centred regular global minimiser
in the chart of Lemmas~\ref{lem:selection}--\ref{lem:variation}, after
reducing the charge interval. The selected state is therefore its
actual positive normalised ground state. Global minimality implies
stationarity along the exactly volume-preserving chart variations
constructed in Lemma~\ref{lem:variation}; that lemma gives constant
negative physical normal flux. This is precisely the boundary
condition required by Theorem~\ref{thm:local}.

\begin{proof}[Proof of Theorem~\ref{thm:main}]
Choose $\beta$ from Lemma~\ref{lem:entrance}, and then choose the
neighbourhood radius $\eta>0$ and $q_{\mathrm{loc}}>0$ small enough
for Theorem~\ref{thm:local} and Lemmas~\ref{lem:selection}
and~\ref{lem:variation} at this exponent.  Reduce a positive threshold
$q_0$ so that the source existence theorem, the constant-flux
conclusion above, Lemma~\ref{lem:entrance} with radius $\eta$, and
$q<q_{\mathrm{loc}}$ all hold whenever $0<q<q_0$.
Fix such a $q$.

Every regular global minimising representative, translated to its
barycentre, is now a volume-$V$ domain $D_\rho$ with
$\|\rho\|_{C^{2,\beta}}<\eta$ and constant ground-state normal
flux.  Theorem~\ref{thm:local} implies $\rho=0$.
At least one such minimiser exists by
\cite[Theorem~1.2]{MMR2025}; consequently
\[
 I_q=E_q(B)=I_q^{\mathrm{ae}}.
\]
This proves both global energy inequalities, without restricting the
competitors to bounded, connected, regular, or nearly spherical
sets.

Suppose first that $\Omega$ is measurable, $|\Omega|=V$, and
$E_q^{\mathrm{ae}}(\Omega)=E_q(B)$.
Lemma~\ref{lem:measurable} gives a quasi-open minimising positivity
set $Q$ with $\Omega=Q$ a.e.  Its regular minimising representative is
a translate of $B$ by the argument just given.  Thus $\Omega$ agrees
a.e. with a translate of $B$.
Conversely, a.e. equality with $B+a$ implies
$X(\Omega)=X(B+a)=H^1_0(B+a)$, and hence
$E_q^{\mathrm{ae}}(\Omega)=E_q(B)$.  This proves both directions of
the measurable equality statement.

Now let $\Omega$ be quasi-open, $|\Omega|=V$, and
$E_q(\Omega)=E_q(B)$.  The same positivity-set argument gives
$\Omega=B+a$ a.e.; translate so that $a=0$.
Let $u\geq0$ attain $E_q(\Omega)$.
Because $u=0$ a.e. off $B$, identity
\eqref{eq:global-lipschitz-support} gives $u\in H^1_0(B)$.
Its energy equals $E_q(B)$, so it is a nonnegative normalised ball
ground state and is strictly positive at every point of $B$.
Its quasi-continuous representative vanishes q.e. off $\Omega$.
The interior smooth representative of the ball ground state agrees
with that representative q.e.; therefore
\begin{equation}
 \operatorname{cap}(B\setminus\Omega)=0.
 \label{eq:global-no-slits}
\end{equation}

For the opposite inclusion use the positivity-set characterisation of
quasi-open sets: there is $w\in H^1(\mathbb R^3)$, $w\geq0$, such
that $\Omega=\{\widetilde w>0\}$ q.e.
This characterisation is recalled in
\cite[Section~2.3]{MMR2025}; its global form follows from the bounded
one by exhaustion.  More explicitly, represent
$\Omega\cap B_n$ by nonnegative global $H^1$ functions $w_n$ and
choose $a_n>0$ with
$\sum_n a_n\|w_n\|_{H^1}<\infty$.  The sum
$w=\sum_n a_nw_n$ converges in $H^1$; its increasing partial sums and
q.e. subsequence convergence show that its quasi-continuous positivity
set is the union of those of the $w_n$, up to a set of zero capacity.
This union is $\Omega$ q.e.
Since $\Omega=B$ a.e., $w=0$ a.e. off $B$.
Equation \eqref{eq:global-lipschitz-support} implies $w\in H^1_0(B)$,
so $\widetilde w=0$ q.e. off $B$.  Hence
\begin{equation}
 \operatorname{cap}(\Omega\setminus B)=0.
 \label{eq:global-no-appendages}
\end{equation}
Together, \eqref{eq:global-no-slits} and
\eqref{eq:global-no-appendages} prove that $\Omega$ equals the ball
q.e.  The two arguments are both necessary: a.e. equality alone does
not exclude deleting a null set of positive capacity.
Conversely, q.e. equality with $B+a$ identifies the two
$H^1_0$ spaces and their energies, proving the reverse implication in
the quasi-open equality statement.  This completes the proof.
\end{proof}

\section{Scaling and scope}\label{sec:scope}
For a volume $|B_R|=R^3V$, the normalised dilation $u_R(x)=R^{-3/2}u(x/R)$ gives
\[
 E_q(R\Omega)=R^{-2}E_{qR}(\Omega).
\]
The same identity holds for the measurable-support relaxation. Consequently Theorem~\ref{thm:main} applies at radius $R$ whenever $0<qR<q_0$, with the same translation and representative qualifications. This is a change of scale of the stated theorem.

The proof does not give the largest admissible charge, a quantitative energy deficit away from the ball, or the geometry after spherical optimality fails. It makes no assertion about the quadratic interaction $D(u,u)$ or about attractive charge. Its local statement concerns the physical flux of the normalised ground state on near-ball domains; the unrestricted global conclusion depends on the cited all-minimiser localisation theorem. In particular, local nondegeneracy alone would not exclude a distant competing minimiser.

\begin{thebibliography}{10}
\bibitem{Cavallina2024families}
L.~Cavallina,
\emph{Nondegeneracy implies the existence of parametrized families of free boundaries},
Journal of Differential Equations \textbf{383} (2024), 1--23.
\href{https://doi.org/10.1016/j.jde.2023.11.011}{doi:10.1016/j.jde.2023.11.011}.
Preprint (under the earlier title \emph{How to construct parametrized families of free boundaries near nondegenerate solutions}): \href{https://arxiv.org/abs/2109.12559}{arXiv:2109.12559}.

\bibitem{MMR2026screened}
D.~Mazzoleni, R.~Moraschi and B.~Ruffini,
\emph{Nonexistence for screened Hartree energies},
preprint (2026), \href{https://arxiv.org/abs/2606.16503}{arXiv:2606.16503}.
\bibitem{MMR2025}
D.~Mazzoleni, C.~B.~Muratov and B.~Ruffini,
\emph{An optimal design problem for a charge qubit},
Communications in Partial Differential Equations \textbf{50} (2025), no.~8, 1029--1073.
\href{https://doi.org/10.1080/03605302.2025.2511919}{doi:10.1080/03605302.2025.2511919}.
Author version: \url{https://pagine.dm.unipi.it/muratov/hartree.pdf}.

\bibitem{MazzoleniRuffini2021}
D.~Mazzoleni and B.~Ruffini,
\emph{A spectral shape optimization problem with a nonlocal competing term},
Calculus of Variations and Partial Differential Equations \textbf{60} (2021), article~114.
Preprint: \href{https://arxiv.org/abs/2009.07699}{arXiv:2009.07699}.

\bibitem{PacardSicbaldi2009}
F.~Pacard and P.~Sicbaldi,
\emph{Extremal domains for the first eigenvalue of the Laplace--Beltrami operator},
Annales de l'Institut Fourier \textbf{59} (2009), no.~2, 515--542.
\href{https://doi.org/10.5802/aif.2438}{doi:10.5802/aif.2438}.

\bibitem{Cavallina2024}
L.~Cavallina,
\emph{Why are the solutions to overdetermined problems usually ``as symmetric as possible''?},
Journal of Geometric Analysis \textbf{34} (2024), article~26.
\href{https://doi.org/10.1007/s12220-023-01467-8}{doi:10.1007/s12220-023-01467-8}.
Preprint: \href{https://arxiv.org/abs/2203.12950}{arXiv:2203.12950}.

\bibitem{MMR2026quadratic}
D.~Mazzoleni, R.~Moraschi and B.~Ruffini,
\emph{Existence and regularity in the small-mass regime for a Hartree--Ohta--Kawasaki shape optimization problem},
preprint (2026), \href{https://arxiv.org/abs/2604.13895}{arXiv:2604.13895}.

\bibitem{GPV2012}
V.~Georgiev, F.~Prinari and N.~Visciglia,
\emph{On the radiality of constrained minimizers to the Schr\"odinger--Poisson--Slater energy},
Annales de l'Institut Henri Poincar\'e C, Analyse non lin\'eaire
\textbf{29} (2012), no.~3, 369--376.
\href{https://doi.org/10.1016/j.anihpc.2011.12.001}{doi:10.1016/j.anihpc.2011.12.001}.

\bibitem{LiebLoss2001}
E.~H.~Lieb and M.~Loss,
\emph{Analysis}, second edition, Graduate Studies in Mathematics, vol.~14,
American Mathematical Society, Providence, RI, 2001.

\bibitem{GT2001}
D.~Gilbarg and N.~S.~Trudinger,
\emph{Elliptic Partial Differential Equations of Second Order},
reprint of the 1998 edition, Classics in Mathematics, Springer, Berlin, 2001.
\href{https://doi.org/10.1007/978-3-642-61798-0}{doi:10.1007/978-3-642-61798-0}.
\end{thebibliography}
\end{document}
